\documentclass[11pt]{article}

\usepackage[a4paper,margin=1in]{geometry}
\usepackage[T1]{fontenc}
\usepackage[utf8]{inputenc}
\usepackage{enumitem}
\usepackage{hyperref}

\title{Paper Notes}
\date{}

\begin{document}
\maketitle

\section*{Current Paper Scope}

The paper should remain simple, elegant, and narrow.

\begin{itemize}[leftmargin=*]
    \item Alternation: transitive active/passive only.
    \item Main conditions: \texttt{CORE}, \texttt{ANOMALOUS}, and \texttt{jabberwocky}.
    \item Main model for the first pass: \texttt{gpt2-large}.
    \item Main experiment only for now: fixed-target structural priming.
\end{itemize}

Things to exclude from the first paper:

\begin{itemize}[leftmargin=*]
    \item PCA
    \item steering
    \item dative materials
    \item clinical framing
    \item semantic-similarity conditions in the first pass
\end{itemize}

\section*{Main Research Question}

Does structural priming in autoregressive language models survive when lexical and
semantic support are progressively depleted, and if so, where in the target
sentence does the residual effect emerge?

\section*{Main Claim to Aim For}

Safe version:

\begin{quote}
Large language models retain measurable structural priming under severe lexical
and semantic degradation, but the effect is attenuated and concentrated in the
structure-diagnostic region of the target.
\end{quote}

Avoid stronger claims such as:

\begin{itemize}[leftmargin=*]
    \item syntax in complete isolation
    \item fully meaning-free stimuli
    \item discovery of a syntax mechanism
\end{itemize}

\section*{Primary Measures}

\begin{itemize}[leftmargin=*]
    \item \texttt{sentence\_pe}: full target sentence priming effect
    \item \texttt{token\_pe}: token-wise priming effect by target position
    \item \texttt{critical\_word\_pe}: full-word priming effect at target word index 2
    \item \texttt{critical\_word\_pe\_mean}: token-normalized full-word priming effect at target word index 2
    \item \texttt{post\_divergence\_pe}: secondary descriptive variable
    \item \texttt{structure\_region\_pe}: legacy coarse suffix metric, useful for diagnostics only
\end{itemize}

Important note:

For active/passive targets, literal token divergence happens very early because
argument order changes immediately. Therefore, \texttt{post\_divergence\_pe} is too
broad to serve as the main structure-sensitive measure. The first pilot also showed
that \texttt{structure\_region\_pe} remains too sensitive to BPE length in the
strict Jabberwocky materials, especially for long nonce verbs. The main analysis
should therefore focus on the word-aligned critical-word metric.

\section*{Current Active Analysis Scripts}

\begin{itemize}[leftmargin=*]
    \item \texttt{scripts/Phase-1/2\_transitive\_token\_priming.py}
    \item \texttt{scripts/Phase-1/3\_summarize\_transitive\_priming.py}
    \item \texttt{scripts/Phase-1/1\_audit\_jabberwocky\_lexicon.py}
    \item \texttt{scripts/Phase-1/2\_audit\_jabberwocky\_semantics.py}
\end{itemize}

The scoring script now writes:

\begin{itemize}[leftmargin=*]
    \item item-level sentence and critical-word scores
    \item token-level target traces
    \item word-level target summaries based on tokenizer-to-word alignment
\end{itemize}

\section*{Why the Jabberwocky Condition Matters}

The Jabberwocky condition must preserve English morphosyntax. That is not a flaw;
it is the point.

What must be preserved:

\begin{itemize}[leftmargin=*]
    \item function words
    \item argument structure
    \item inflectional morphology
    \item passive cues such as \texttt{was} and \texttt{by}
\end{itemize}

What should be minimized:

\begin{itemize}[leftmargin=*]
    \item ordinary lexical semantics of content words
    \item obvious form overlap with existing English words
    \item direct recycling of PrimeLM vocabulary
\end{itemize}

\section*{Jabberwocky Lexicon Revision}

An earlier exploratory Jabberwocky lexicon was archived and replaced by a stricter
nonce lexicon.

Current active files:

\begin{itemize}[leftmargin=*]
    \item \texttt{corpora/transitive/jabberwocky\_transitive.csv}
    \item \texttt{corpora/transitive/jabberwocky\_transitive\_bpe\_filtered.csv}
    \item \texttt{corpora/transitive/jabberwocky\_transitive\_strict\_vocabulary.json}
    \item \texttt{corpora/transitive/jabberwocky\_transitive\_strict\_bpe\_filtered\_vocabulary.json}
    \item \texttt{scripts/corpora\_generators/0\_generate\_jabberwocky\_transitive.py}
\end{itemize}

Archived legacy files:

\begin{itemize}[leftmargin=*]
    \item \texttt{corpora/archive/legacy/transitive/jabberwocky\_transitive.csv}
    \item \texttt{corpora/archive/legacy/transitive/jabberwocky\_transitive\_legacy\_vocabulary.json}
    \item \texttt{scripts/archive/legacy/corpora\_generators/0\_generate\_jabberwocky\_transitive.py}
\end{itemize}

\section*{Lexical Audit Note for the Paper}

This should be reported briefly in Methods or Stimuli Validation.

Suggested content:

\begin{quote}
We replaced an earlier exploratory Jabberwocky lexicon with a stricter nonce
lexicon preserving English morphosyntax while minimizing lexical overlap with
English and PrimeLM vocabulary. A lexical audit of the revised lexicon found no
exact English matches, no PrimeLM vocabulary overlaps, no stem overlaps, and no
neighbors within edit distance $\leq 2$ under our screening criteria.
\end{quote}

The detailed audit tables should go in the appendix or supplement, not the main
results section.

\section*{Semantic Leakage Audit Note}

We also ran a model-based representational audit using GPT-style lexical and
contextual representations.

Interpretation to remember:

\begin{itemize}[leftmargin=*]
    \item Low lexical similarity is evidence against direct lexical anchoring.
    \item Contextual similarity in neutral noun/verb frames is expected to reflect
    grammatical slot information, not just meaning.
    \item Therefore, contextual similarity does not by itself prove semantic
    contamination.
\end{itemize}

Main conclusion from the comparison:

\begin{itemize}[leftmargin=*]
    \item The new strict lexicon reduced lexical anchoring relative to the legacy
    Jabberwocky lexicon.
    \item The improvement is strongest for nouns and past-tense verb forms.
    \item Present-tense verb forms remain moderately close to English verbs at the
    lexical level, which is expected because English morphology is intentionally
    preserved.
\end{itemize}

Suggested wording:

\begin{quote}
A model-based audit further indicated reduced lexical anchoring in the revised
Jabberwocky lexicon relative to the earlier exploratory version, especially for
noun forms and past-tense verb forms.
\end{quote}

\section*{Tokenizer Handling and How to Justify It}

For \texttt{gpt2-large}, there are no ordinary unknown tokens. GPT-2 uses
byte-level BPE, so a nonce form is decomposed into valid subword pieces rather
than mapped to an \texttt{<unk>} symbol.

This means the relevant confound is not unknown-token failure but segmentation
length:

\begin{itemize}[leftmargin=*]
    \item within a single priming comparison, the target sentence is identical
    under the congruent and incongruent prime, so target tokenization is held
    constant
    \item across items and structures, longer BPE decompositions can distort
    summed token-level effects
\end{itemize}

Therefore:

\begin{itemize}[leftmargin=*]
    \item we should not justify the method in terms of ``handling unknown tokens''
    \item we should justify it in terms of controlling subword segmentation length
    \item the main safeguard is to compute priming over the full critical word span
    rather than a single BPE token or a long suffix
    \item the token-normalized form \texttt{critical\_word\_pe\_mean} is especially
    important for Jabberwocky because active nonce verbs often decompose into
    4--6 GPT-2 tokens whereas passive auxiliaries are usually 1 token
\end{itemize}

Suggested wording:

\begin{quote}
Because GPT-2 uses byte-level BPE, nonce items are not mapped to unknown symbols
but decomposed into subword sequences. We therefore treat subword segmentation
length, rather than unknown-token failure, as the relevant tokenization confound.
To reduce this confound, priming at the critical word is computed over the full
word span and reported both as a summed effect and as a token-normalized mean.
\end{quote}

\section*{Possible Additional Tokenization Control}

This control step has now been implemented for the transitive Jabberwocky
materials.

Current filter specification:

\begin{itemize}[leftmargin=*]
    \item noun forms kept only if GPT-2 tokenization length is 4 or 5 BPEs
    \item verb stems kept only if both present and past inflected forms are 4 or 5 BPEs
\end{itemize}

Current filtered outputs:

\begin{itemize}[leftmargin=*]
    \item \texttt{behavioral\_results/jabberwocky\_tokenizer\_length\_audit/tokenizer\_length\_summary.csv}
    \item \texttt{behavioral\_results/jabberwocky\_tokenizer\_length\_audit/noun\_token\_lengths.csv}
    \item \texttt{behavioral\_results/jabberwocky\_tokenizer\_length\_audit/verb\_token\_lengths.csv}
    \item \texttt{corpora/transitive/jabberwocky\_transitive\_strict\_bpe\_filtered\_vocabulary.json}
    \item \texttt{corpora/transitive/jabberwocky\_transitive\_bpe\_filtered.csv}
\end{itemize}

Current retained inventory under the 4--5 BPE filter:

\begin{itemize}[leftmargin=*]
    \item nouns: 80 $\rightarrow$ 65
    \item verb stems: 40 $\rightarrow$ 34
\end{itemize}

This means the stricter BPE control remains feasible without collapsing the
vocabulary.

If reviewers still push on subword length, one stronger control remains available:

\begin{itemize}[leftmargin=*]
    \item tighten the allowed token-count band further
    \item explicitly balance token-count distributions across present and past forms
    \item optionally match noun token counts even more aggressively
\end{itemize}

This is not necessary before the first full run, because the current paper path
already uses the BPE-filtered Jabberwocky corpus.

\section*{What to Report vs.\ What to Avoid}

Report:

\begin{itemize}[leftmargin=*]
    \item that the Jabberwocky lexicon was revised
    \item that lexical-confound screening was performed
    \item that the revised lexicon is cleaner than the legacy one
    \item that English morphosyntax was intentionally preserved
\end{itemize}

Avoid:

\begin{itemize}[leftmargin=*]
    \item claiming that semantics has been eliminated absolutely
    \item claiming that the stimuli are meaning-free in a philosophical sense
    \item turning the paper into a lexical audit paper
\end{itemize}

\section*{Best Framing}

Best framing for the Jabberwocky condition:

\begin{quote}
lexically content-degraded while preserving English morphosyntax
\end{quote}

Good alternatives:

\begin{itemize}[leftmargin=*]
    \item Jabberwocky-style nonce lexicalization
    \item content-degraded morphosyntactic materials
\end{itemize}

\section*{Immediate Next Step}

Pilot reruns with the revised scoring pipeline are now complete. The main point
to remember is:

\begin{itemize}[leftmargin=*]
    \item the coarse suffix metrics still look distorted in strict Jabberwocky
    \item the word-aligned critical-word metric is much cleaner
    \item Jabberwocky active priming is above zero at the critical word, even
    though sentence-level and suffix-level summaries can remain negative
\end{itemize}

Pilot values to remember from the revised pipeline (\texttt{max-items = 200}):

\begin{itemize}[leftmargin=*]
    \item \texttt{CORE} active critical-word PE: about 0.057
    \item \texttt{ANOMALOUS} active critical-word PE: about 0.0004
    \item strict \texttt{jabberwocky} active critical-word PE: about 0.122
    \item \texttt{CORE} passive critical-word PE: about 0.227
    \item \texttt{ANOMALOUS} passive critical-word PE: about 0.151
    \item strict \texttt{jabberwocky} passive critical-word PE: about 0.922
\end{itemize}

Therefore the next real step is the full Experiment 1 run with the active strict
Jabberwocky corpus:

\begin{itemize}[leftmargin=*]
    \item \texttt{CORE}
    \item \texttt{ANOMALOUS}
    \item new strict BPE-filtered \texttt{jabberwocky}
\end{itemize}

The paper is viable only if:

\begin{itemize}[leftmargin=*]
    \item the token-level curves are interpretable
    \item \texttt{critical\_word\_pe} or \texttt{critical\_word\_pe\_mean} remains
    above zero for strict Jabberwocky
    \item the effect ordering remains interpretable across \texttt{CORE},
    \texttt{ANOMALOUS}, and strict \texttt{jabberwocky}
    \item active/passive asymmetry is discussed as a property of the paradigm,
    not overinterpreted as a direct structural comparison
\end{itemize}

\section*{Analysis Methods Update (23 March 2026)}

This section supersedes the earlier pilot-only status notes.

\subsection*{Data and run configuration}

\begin{itemize}[leftmargin=*]
    \item Full Experiment 1 was run with \texttt{gpt2-large} using
    \texttt{make transitive-priming}.
    \item Active preset: \texttt{paper\_main}.
    \item Conditions: \texttt{CORE}, \texttt{ANOMALOUS}, strict BPE-filtered
    \texttt{jabberwocky}.
    \item Corpus paths are now \texttt{corpora/transitive/*} in the active preset.
    \item Checkpoint/resume is enabled per condition to avoid rerun loss.
\end{itemize}

\subsection*{Primary scoring outputs}

\begin{itemize}[leftmargin=*]
    \item Item-level file: 90,000 rows total (30,000 per condition; active and passive targets per item).
    \item Token-level file: 951,023 rows total.
    \item Word-level file: 630,000 rows total.
    \item Summary files and report:
    \texttt{behavioral\_results/transitive\_token\_profiles/*}.
\end{itemize}

\subsection*{Statistical analysis implementation}

\begin{itemize}[leftmargin=*]
    \item New script:
    \texttt{scripts/Phase-1/5\_analyze\_transitive\_statistics.py}.
    \item New make target: \texttt{make transitive-stats}.
    \item Primary dependent variable: \texttt{sentence\_pe\_mean}
    (passive minus active, item-paired).
    \item Secondary dependent variables:
    \texttt{sentence\_pe}, \texttt{post\_divergence\_pe\_mean},
    \texttt{critical\_word\_pe\_mean}.
    \item Item-paired inference:
    one-sample paired contrast with
    \begin{itemize}[leftmargin=*]
        \item sign-flip permutation test (\(n=10{,}000\))
        \item bootstrap 95\% CI (\(n=10{,}000\))
        \item one-sample \(t\)-test and paired \(d_z\)
    \end{itemize}
    \item Confound-aware model on paired deltas:
    \texttt{delta\_sentence\_pe\_mean \(\sim\) condition + z\_delta\_target\_length + z\_delta\_critical\_word\_token\_count}
    with HC3 robust SE.
    \item Per-condition robustness model:
    random-intercept LMM by item,
    \texttt{sentence\_pe\_mean \(\sim\) target\_structure + z\_target\_length + z\_critical\_word\_token\_count}.
\end{itemize}

\subsection*{Key results from full run}

\paragraph{Primary metric (\texttt{sentence\_pe\_mean}, passive minus active).}

\begin{itemize}[leftmargin=*]
    \item \texttt{CORE}: \(+0.1434\), 95\% CI \([0.1375,\,0.1493]\), \(d_z=0.389\).
    \item \texttt{ANOMALOUS}: \(+0.2500\), 95\% CI \([0.2433,\,0.2566]\), \(d_z=0.594\).
    \item \texttt{jabberwocky}: \(+0.4975\), 95\% CI \([0.4921,\,0.5029]\), \(d_z=1.482\).
    \item All permutation \(p\)-values are at the floor (\(0.0001\) with 10k permutations).
\end{itemize}

\paragraph{Raw sentence-sum metric (\texttt{sentence\_pe}).}

\begin{itemize}[leftmargin=*]
    \item \texttt{CORE}: \(+1.208\)
    \item \texttt{ANOMALOUS}: \(+1.893\)
    \item \texttt{jabberwocky}: \(+8.994\)
\end{itemize}

This raw scale is strongly length-sensitive and should remain secondary.

\paragraph{Confound-aware delta regression.}

\begin{itemize}[leftmargin=*]
    \item \texttt{z\_delta\_target\_length}: not significant (\(p=0.88\)).
    \item \texttt{z\_delta\_critical\_word\_token\_count}: positive and significant
    (\(\beta=0.0884,\ p\approx 1.58\times 10^{-19}\)).
    \item \texttt{jabberwocky} remains strongly above baseline after covariates.
\end{itemize}

\paragraph{Per-condition LMM robustness (passive coefficient).}

\begin{itemize}[leftmargin=*]
    \item \texttt{CORE}: \(0.1704,\ p\approx 2.96\times 10^{-10}\).
    \item \texttt{ANOMALOUS}: \(0.4098,\ p\approx 4.01\times 10^{-42}\).
    \item \texttt{jabberwocky}: \(0.6497,\ p\approx 0\) (numerically underflowed).
    \item Final LMM fits converged for all three conditions after optimizer fallback.
\end{itemize}

\subsection*{Current interpretation for writeup}

\begin{itemize}[leftmargin=*]
    \item The model shows a robust passive-over-active priming asymmetry in all conditions.
    \item The asymmetry remains in strict Jabberwocky under token-normalized and covariate-controlled analyses.
    \item The largest effects are concentrated in morphosyntactic cue zones
    (especially passive auxiliaries and function words), not in nonce lexical stems.
    \item Therefore the safe claim is morphosyntactic-cue-sensitive priming under
    lexical content degradation, not fully isolated ``pure syntax''.
\end{itemize}

\subsection*{Files to cite in Methods/Results drafting}

\begin{itemize}[leftmargin=*]
    \item \texttt{behavioral\_results/transitive\_token\_profiles/transitive\_item\_summary.csv}
    \item \texttt{behavioral\_results/transitive\_token\_profiles/transitive\_report.md}
    \item \texttt{behavioral\_results/transitive\_token\_profiles/stats/paired\_effects.csv}
    \item \texttt{behavioral\_results/transitive\_token\_profiles/stats/delta\_regression\_coefficients.csv}
    \item \texttt{behavioral\_results/transitive\_token\_profiles/stats/lmm\_condition\_coefficients.csv}
    \item \texttt{behavioral\_results/transitive\_token\_profiles/stats/analysis\_report.md}
\end{itemize}

\end{document}
